% !Mode:: "TeX:UTF-8"

\chapter{绪\hspace{1em}论}
\section{课题背景及研究的目的和意义}
\subsection{课题背景}

近年来，大语言模型（Large Language Models, LLMs）在模型规模、预训练数据、指令微调以及人类反馈强化学习等方面取得了显著进展，
展现出较强的语言理解、知识整合与复杂推理能力。随着模型能力的不断提升，大语言模型已逐渐从通用文本生成工具发展为面向复杂任务的智能推理系统，
并在科学研究、智能教育、代码生成、数学求解、知识问答和决策辅助等关键领域中承担越来越重要的角色。尤其是在需要多步骤推理的场景中，
大语言模型通常能够生成较为完整的中间推理过程，为问题求解提供了新的技术路径和应用可能。

然而，大语言模型在复杂推理任务中的可靠性仍然面临严峻挑战。尽管模型能够生成形式上连贯、逻辑上看似合理的推理链，
但其中常常包含隐蔽的步骤错误，例如前提理解偏差、事实引用错误、数学计算失误、逻辑跳跃、因果关系混淆以及结论与中间推导不一致等问题。
此类错误往往不易被用户及时发现，并可能随着推理链的推进不断累积，最终导致错误结论。与简单事实性错误相比，推理过程中的步骤错误更具隐蔽性和危害性，
因为模型生成的解释性文本可能增强用户对错误答案的信任，从而在实际应用中引发误导性判断甚至严重后果。

在科学研究和教育等高可信度需求场景中，模型推理错误的影响尤为突出。例如，在科学问题求解中，错误的推理步骤可能导致错误假设或无效结论；在教育场景中，
模型若给出看似合理但实际错误的解题过程，可能会误导学习者形成错误认知；在决策辅助任务中，推理链中的局部错误也可能放大为整体决策风险。
因此，系统性地识别、分析并缓解大语言模型推理过程中的步骤错误，对于提升模型输出的可信度、可解释性和安全性具有重要意义。

目前，国内外学者已围绕大语言模型的错误生成、幻觉现象、事实一致性、推理能力评估以及模型对齐等方向开展了大量研究\cite{ref01_huang2025survey,ref02_che2023nlp}，
并提出了包括提示优化、检索增强、结果验证、自一致性采样和外部工具辅助等多种改进方法。
这些研究在一定程度上推动了大语言模型可靠性问题的解决。然而，现有工作大多集中于最终答案的正确性评估或特定类型错误的检测，
对推理过程中“错误如何发生以及如何针对性缓解”等问题仍缺乏系统性研究。
换言之，当前仍缺少一套贯穿错误识别、错误归因、错误分析与错误缓解的完整方法框架。
这一不足限制了研究者对大语言模型推理缺陷的深入理解，也阻碍了大语言模型在复杂高风险场景中的进一步应用。

\subsection{研究的目的和意义}

本文面向大语言模型在复杂推理过程中的步骤错误问题，旨在构建一套错例分析与缓解流程。
具体而言，本文将围绕推理步骤错误的识别、分类、成因分析与优化缓解展开研究：
首先，建立面向大语言模型推理过程的错误识别标准与分析流程，对模型输出中的中间推理步骤进行细粒度检查；
其次，结合错误分析结果，探索针对性的缓解方法，以减少模型在推理任务中的错误发生率，提高模型输出的正确性与稳定性。
最后，结合错误检测以及使用增强手段后的模型，构建协同进行的可信回答生成流程。

本研究的意义主要体现在理论价值和应用价值两个方面。从理论层面来看，本文试图从“过程错误”而非仅从“最终答案错误”的角度审视大语言模型的推理能力，
有助于更加细致地揭示模型在复杂任务中的能力边界与缺陷模式。通过构建错误识别、错误归因和错误缓解相结合的系统性流程，
本文预期能够为大语言模型错误诊断提供一种更加完整的方法论参考，并进一步完善大语言模型可靠性分析与推理能力评估的理论框架。

从应用层面来看，本文的研究有助于提升大语言模型在科学、教育和决策辅助等场景中的可信使用水平。通过对错例进行细粒度分析，
可以帮助研究者和开发者更准确地定位模型失败的原因，从而避免仅依赖最终答案正确率进行粗粒度评估。同时，针对不同错误类型设计相应的缓解策略，
也有助于提升模型在复杂推理任务中的鲁棒性、可解释性和安全性。对于实际部署而言，该研究能够为大语言模型的评测、优化和应用提供更具操作性的参考，
降低模型错误推理带来的潜在风险。

开展大语言模型错例成因与缓解分析研究，有助于深化对大语言模型推理机制和局限性的认识，为构建更加可靠、
可控和可信的大语言模型应用系统提供重要支撑。本文预期提出一套面向大语言模型推理缺陷的辨别、分析与缓解方法，
为后续相关研究和工程实践提供可复用的分析框架与技术路径。

\section{国内外在该方向的研究现状及分析}
\subsection{大语言模型错误识别}
语言模型能够生成类人的响应,这使得检测这些错误尤为困难。人工检查依赖专家
经验,这种模式人力成本高昂且难以规模化。至于自动化方法，则通常采用静态和动态两
类基准测试。基准测试主要为了对模型进行评估和排行，对识别模型的错误作用有限。
红队测试被用于检测大语言模型错误\cite{ref03_dinan2019build}，但人工测试成本高昂。基于语言模型的红队
测试因而被提出\cite{ref04_perez2022red}。另外，Cheng 等人\cite{ref05_cheng2024autodetect}使用大语言模型多智能体框架检测大语言模
型错误。
对于大语言模型事实层面错误检查有丰富的研究，可分为基于事实核查的检测，以及
基于不确定性估计的检测。Min 等人\cite{ref06_min2023factscore}设计了专为评估长文本生成而设计的细粒度事实度
量方法。该方法首先将生成内容分解为原子事实,随后计算其受可靠知识源支持的比例。
Chern 等人\cite{ref07_chern2023factool}提出了一个统一框架,通过利用一系列专门用于证据收集的外部工具,使 LLM
具备识别事实错误的能力。Dhuliawala 等人\cite{ref08_dhuliawala2023cove}提出 CoVe 方法,让 LLM 先生成针对草稿回
答的验证问题并评估答案与原始回复的一致性。Kadavath 等人\cite{ref09_kadavath2022language}和 Zhang 等人\cite{ref10_zhang2024selfalignment}通过模型内部
知识计算概率 \(p(\mathrm{True})\) 来评估布尔问题的回答事实性。Varshney 等人\cite{ref11_varshney2023stitch}通过计算关键概念
内的最小词元概率来量化模型对这些概念的不确定性。Manakul 等人\cite{ref12_manakul2023selfcheckgpt}通过从 LLM 中对同
一提示采样多个响应,通过评估事实陈述的一致性来检测幻觉。
大语言模型在推理过程中应该忠实于上下文和用户指令。在大语言模型忠实层面错误
检查方面，存在着丰富的研究。如基于衡量生成内容与源内容之间关键事实的重叠程度、
利用从相关任务数据训练的分类器以及基于回答的指标\cite{ref13_falke2019ranking}。不确定性不仅在事实性层面上
可以用于检测，也能在忠实性层面进行检测。Lakshminarayanan 等人\cite{ref14_lakshminarayanan2017simple}提出的深度集成方法
为预测不确定性估计提供了常用技术基础。
随着 deepseek-reasoner 等能生成长链思维推理步骤来增强推理能力的 o1 类模型大语
言模型出现，思维链中错误的检测也是引起了广泛关注。He 等人\cite{ref15_he2025large}构建了 DeltaBench 基
准数据集,其中包含来自不同 o1 类模型针对各类推理任务生成的长链思维,以衡量模型在
长链思维推理中检测错误的能力，并在该基准上测试了过程奖励模型与大语言模型在思维
链错误识别上的能力。You 等人\cite{ref16_you2025probabilistic}采取将上下文与思维链分解为基础声明，使用大语言模
型对基础声明评价合理性概率并蒙特卡洛采样对思维链条采样以计算各条声明蕴含稳定
性。DeepSeekMath-V2\cite{shao2025deepseekmath}收集大语言模型在数学题上错误的样本并标注。
训练专门的大语言模型对数学证明中是否存在错误进行分析判断，并给出可信评分。进而对模型推理证明过程
中错误进行检测。



\subsection{大语言模型错误的成因分析}
面向推理过程，大语言模型错误的成因大致可分为四类：非完美解码策略、Softmax
瓶颈、过度自信、推理失败。
目标为优化高概率输出的解码策略(如束搜索)会导致文本质量严重退化、陷入。束搜
与贪婪解码在多样性和表现性能方面低于采样策略\cite{ref17_holtzman2020curious}。解码策略与模型内部表示也会影响生成事实性，
例如 DoLa 通过对比不同层表示来改善事实性\cite{ref18_chuang2023dola}；在检索增强场景中，归因性与流畅性之间也存在权衡\cite{ref19_aksitov2023characterizing}。
采样温度的升高会导致词元概率分布更加均匀,从而增加了采样不常出现词元的概率，加剧了幻觉的风险。
过度自信问题根源在于过度关注部分生成内容,牺牲对上下文的忠实遵循。语言模型的
注意力机制往往表现出局部聚焦特性,会优先关注邻近词汇,导致对上下文的关注度明显不足\cite{ref20_shi2023trusting}。
使用 softmax 层结合词嵌入计算词预测相关的输出概率分布限制了语言模型输出所需
的分布。Chang 和 McCallum\cite{ref21_chang2022softmax}发现,当输出词嵌入空间中的期望分布呈现多模态时,语言模
型难以准确地 将所有模态中的词优先排序为下一个最可能的词,这也引入了幻觉的风险。
知识的有效利用与推理能力密不可分。即使大语言模型拥有必要的知识,由于其推理
能力的局限性\cite{ref22_zheng2023why},可能难以产生准确的结果。尤其是问题之间存在多重关联时推理能力的
局限性容易显现。另外，大语言模型在推理过程中无用或错误的推理步骤也会影响到后续
的推理过程，而造成错误的推理步骤的一个重要的原因是错误地应用知识。



\subsection{大语言模型错误的优化改进}
推理相关错误的缓解与优化可以解码相关和数据相关两个角度考虑。其中解码相关的
缓解错误手段可分为事实性增强解码与忠实性增强解码解码两方面。
事实性解码可在初始解码过程中修改概率分布。Lee 等人\cite{ref23_lee2022factuality}提出了事实核心采样算法,
该算法在句子生成过程中动态调整核心概率 p，在生成事实内容和保持输出多样性之间取
得了平衡。Chuang 等人\cite{ref18_chuang2023dola}利用 Transformer 事实知识的分层编码特性,注意到较低层级的
信息捕获于早期层,而语义信息则体现在后期层，通过强调高层知识并弱化低层知识,使大
语言模型更符合事实、从而减少幻觉的潜力。后编辑解码则旨在利用大型语言模型的自我
校正能力来优化原始生成内容,而无需依赖外部知识库。Dhuliawala 等人\cite{ref08_dhuliawala2023cove}提出了验证链
(COVE)。该方法从初始草稿开始,首先制定验证问题,然后系统回答这些问题,最终生成改
进后的修订响应。
忠实性增强旨在增强上下文一致性与逻辑一致性。Shi 等人\cite{ref20_shi2023trusting}提出了上下文感知解码
(CAD),该方法通过对比性框架修改模型的原始输出分布。通过放大有上下文与无上下文时
输出概率的差异,使大语言模型更关注上下文信息。逻辑推理严谨的逻辑计算,而纯文本缺
乏精确的逻辑结构,容易导致不忠实的推理。Xu 等人\cite{ref27_xu2024faithful}提出了 SymbCoT,在思维链中融入符
号表达式来描述中间推理步骤。DeepSeekMath-V2\cite{shao2025deepseekmath}训练专门的大语言模型对数学证明中是否存在错误进行分析判断，并给出可信评分。进而对模型推理证明过程
中错误进行检测，将检测信号用于构建强化学习奖励函数并在拥有推理验证能力的模型基础上进行数学证明生成强化学习，以优化中间推理步骤的忠实性
另外，有方法采用混合 Softmax，使用多个隐藏状态多次计算 softmax 并合并结果分
布\cite{ref28_yang2019mixtape}，也有方法引入指针网络使 LLM 能复制上下文词汇\cite{ref29_chang2023revisiting},旨在突破 softmax 瓶颈，减少
上下文幻觉。
数据相关的常见缓解错误手段为检索增强生成，能有效缓解因知识确实导致的错误，
一定程度上进行知识与大语言模型的解耦。根据检索时机不同,检索增强生成分为三类:
一次性检索 、迭代检索与事后检索。一次性检索旨在将从单次检索中获取的外部知识直
接附加到 LLMs 的提示前。Ram 等人\cite{ref30_ram2023incontext}提出了上下文检索增强语言模型(In-context RALM),
该模型采用了一种 简单而有效的策略,即将选定的文档附加到 LLMs 的输入文本前。
当面临多步推理和长文本问答等复杂挑战,允许在生成过程中持续收集知识的迭代检
索方法有更出色的表现。Press 等人\cite{ref31_press2022measuring}提出了自我提问法。自我提问法会明确每一步待解
决的问题,随后基于后续问题加入搜索动作，在每个推理步骤中融入外部知识,并基于持续
进行的推理进一步指导检索过程,从而减少推理链中的事实错误。
事后检索方法,通过基于检索的后续修订来优化大语言模型的输出。Zhao 等人\cite{ref32_zhao2023verify}提出
验证和编辑框架，针对一致性低于平均水平的推理链生成验证性问题,并基于检索到的知
识优化推理依据,从而确保回答更具事实性。




\section{本文的主要研究内容}
本文围绕大语言模型在复杂问答生成过程中的错误识别、错误缓解和可信回答生成展开研究。
与仅评价最终答案是否正确的工作不同，本文关注模型在“思考-检索-信息生成-回答”过程中的知识调用和推理一致性问题。
研究目标是先建立能够识别过程错误的验证方法，再将错误分析结果转化为可学习的反思监督信号，
最后在生成模型与验证模型协同工作的框架下提升问答结果的可信程度。



全文结构安排如下。
第一章介绍课题背景、研究意义和国内外研究现状，并给出本文的研究内容与章节组织。
第二章研究大语言模型生成过程的错误识别问题，构建评测数据，比较 PRM 异常检测与 LLM-as-a-judge 方法，
并确定后续实验采用的反思验证模型。针对大语言模型在自检索问答中可能出现的问题理解偏差、检索查询偏移、中间信息错误和答案整合不一致等现象，
本文构建基于 SSRL rollout 的过程错误评测数据，并比较过程奖励模型与大语言模型评判器两类错误识别方法。
其中，过程奖励模型方法利用片段级奖励分数和 Z-score 异常检测识别低质量推理片段；
大语言模型评判器方法则根据问题、完整生成轨迹和标准答案输出错误判断、错误说明和反思反馈。
实验结果表明，基于评判模型的错误识别更适合处理需要全局语义比较的过程错误，
也能为后续错误缓解提供可直接使用的文本监督。

第三章研究面向错误缓解的策略，将第二章得到的错误分析和反思反馈转化为监督微调数据，
并验证反思增强初始化对自检索强化学习的影响。针对结果奖励信号稀疏、难以定位中间推理错误的问题，
本文提出一种先进行反思验证监督微调、再继续执行自检索强化学习的错误缓解策略。
该策略使用强验证模型对模型生成轨迹进行过程级标注，生成包含错误原因和修正建议的反思验证数据，
并以此训练反思增强模型，使其在强化学习前获得更明确的错误感知和修正先验。
随后，本文在多个开放域问答和多跳问答数据集上比较该方法与直接 SSRL 训练的差异，
验证过程级反思监督对模型最终回答准确率和训练过程中最佳可达性能的影响。


第四章研究多模型协同可信问答方法，在生成模型与验证模型之间引入迭代反馈机制，
实现带有过程质量控制的候选答案生成与筛选。在错误识别和错误缓解的基础上，本文进一步构建由策略模型、标准答案匹配模块和验证模型组成的迭代式可信问答流程。
策略模型负责生成多个候选自检索回答，标准答案匹配模块用于判断候选答案是否命中标准答案，
验证模型则从推理过程、知识调用和答案一致性角度筛选可信轨迹。
对于未通过验证的样本，系统将验证反馈重新输入策略模型，引导其在后续轮次中进行反思式再生成。
该流程将生成、验证和反馈整合为闭环，能够在提高答案命中率的同时过滤过程不可靠的候选结果。





% Local Variables:
% TeX-master: "../thesis"
% TeX-engine: xetex
% End:
